\documentclass[12pt,a4paper]{article}

% ==================== GÓI CƠ BẢN ====================
\usepackage[utf8]{inputenc}
\usepackage[T5]{fontenc}
\usepackage[vietnamese]{babel}
\usepackage{geometry}
\usepackage{parskip}
\usepackage{enumitem}
\usepackage{longtable}
\usepackage{array}
\usepackage{xcolor}
\usepackage{hyperref}
\usepackage{fancyhdr}
\usepackage{lastpage}
\usepackage{amsmath}

% ==================== ĐỊNH DẠNG TRANG ====================
\geometry{
    top=2cm,
    bottom=2cm,
    left=2.5cm,
    right=2cm
}

\setlength{\parindent}{0pt}
\setlength{\parskip}{0.7em}
\setlist[itemize]{leftmargin=1.8em, itemsep=0.35em}
\setlist[enumerate]{leftmargin=1.8em, itemsep=0.35em}

\definecolor{ailmsblue}{RGB}{25,84,166}
\definecolor{ailmsgray}{RGB}{90,90,90}
\definecolor{ailmsred}{RGB}{180,40,40}

\hypersetup{
    colorlinks=true,
    linkcolor=ailmsblue,
    urlcolor=ailmsblue,
    citecolor=ailmsblue,
    pdftitle={Chính sách Thương mại và Thanh toán AILMS},
    pdfauthor={Phòng Tài chính và Kinh doanh AILMS}
}

% ==================== HEADER / FOOTER ====================
\pagestyle{fancy}
\fancyhf{}
\fancyhead[L]{\textcolor{ailmsgray}{AILMS}}
\fancyhead[R]{\textcolor{ailmsgray}{Chính sách Thương mại \& Thanh toán}}
\fancyfoot[C]{Trang \thepage/\pageref{LastPage}}
\renewcommand{\headrulewidth}{0.4pt}
\renewcommand{\footrulewidth}{0.4pt}

% ==================== THÔNG TIN TÀI LIỆU ====================
\title{
    \textbf{\LARGE CHÍNH SÁCH THƯƠNG MẠI\\
    VÀ THANH TOÁN AILMS}
}
\author{Phòng Tài chính và Kinh doanh AILMS}
\date{Phiên bản 1.0 -- Ngày hiệu lực: \today}

\begin{document}

\maketitle

\begin{center}
    \textit{Tài liệu này quy định nguyên tắc niêm yết giá, khuyến mại,
    đặt hàng, thanh toán, hoàn tiền, đối soát và chia sẻ doanh thu
    đối với các dịch vụ đào tạo trên hệ thống AILMS.}
\end{center}

\vspace{0.8cm}

\begin{longtable}{|>{\bfseries}p{4cm}|p{9.5cm}|}
    \hline
    Tên tài liệu &
    Chính sách Thương mại và Thanh toán AILMS \\
    \hline
    Phiên bản & 1.0 \\
    \hline
    Đơn vị quản lý &
    Phòng Tài chính và Kinh doanh AILMS \\
    \hline
    Đối tượng áp dụng &
    Khách hàng, học viên, người giám hộ, giảng viên,
    bộ phận tài chính, bộ phận kinh doanh và quản trị viên liên quan \\
    \hline
    Phạm vi áp dụng &
    Giỏ hàng, đơn hàng, mã giảm giá, thanh toán, hoàn tiền,
    hóa đơn, đối soát và chia sẻ doanh thu \\
    \hline
    Ngày hiệu lực & \today \\
    \hline
\end{longtable}

\tableofcontents
\newpage

% =========================================================
\section{Mục đích và phạm vi áp dụng}

Chính sách này nhằm:

\begin{itemize}
    \item Công khai nguyên tắc xác định giá và khuyến mại.
    \item Quy định quy trình tạo và xử lý đơn hàng.
    \item Bảo đảm giao dịch được xác minh an toàn.
    \item Ngăn ngừa việc cấp quyền học khi chưa thanh toán thành công.
    \item Quy định nguyên tắc hoàn tiền và xử lý tranh chấp.
    \item Thiết lập quy trình đối soát và chia sẻ doanh thu.
    \item Bảo vệ quyền lợi hợp pháp của khách hàng và các bên cung cấp dịch vụ.
\end{itemize}

Chính sách áp dụng đối với các giao dịch mua khóa học,
gói đào tạo và dịch vụ liên quan trên AILMS.

Điều kiện riêng được công bố tại từng khóa học hoặc chương trình
được áp dụng cùng với tài liệu này.

% =========================================================
\section{Giải thích thuật ngữ}

\begin{itemize}
    \item \textbf{Sản phẩm:} Khóa học, gói đào tạo,
    buổi tư vấn hoặc dịch vụ được cung cấp trên AILMS.

    \item \textbf{Gói đào tạo:} Phương án cung cấp khóa học theo hình thức
    tự học, lớp nhóm, một-một hoặc hình thức khác.

    \item \textbf{Giỏ hàng:} Danh sách sản phẩm người dùng lựa chọn
    trước khi tạo đơn hàng.

    \item \textbf{Đơn hàng -- Order:} Bản ghi giao dịch thể hiện
    sản phẩm, giá, khuyến mại, số tiền và trạng thái thanh toán.

    \item \textbf{Order Item:} Một sản phẩm hoặc gói đào tạo
    được ghi nhận trong đơn hàng.

    \item \textbf{Coupon:} Mã được sử dụng để áp dụng ưu đãi
    theo điều kiện đã cấu hình.

    \item \textbf{Payment:} Bản ghi thể hiện một lần hoặc một quá trình
    thanh toán gắn với đơn hàng.

    \item \textbf{Approval:} Hành động người mua đồng ý thanh toán
    trên trang của nhà cung cấp thanh toán.

    \item \textbf{Capture:} Bước Backend yêu cầu nhà cung cấp
    hoàn tất và xác nhận khoản thanh toán.

    \item \textbf{Refund:} Giao dịch hoàn trả toàn bộ hoặc một phần
    số tiền đã thanh toán.

    \item \textbf{Chargeback:} Quy trình chủ tài khoản thanh toán
    yêu cầu tổ chức phát hành hoặc nhà cung cấp thanh toán
    xem xét và đảo ngược giao dịch.

    \item \textbf{Revenue Share:} Cơ chế phân chia doanh thu
    giữa AILMS và giảng viên hoặc đối tác.

    \item \textbf{Giá trị thực thu:} Khoản tiền được xác nhận thu thành công
    sau các khoản giảm giá và điều chỉnh được quy định.
\end{itemize}

% =========================================================
\section{Nguyên tắc thương mại}

AILMS áp dụng các nguyên tắc:

\begin{itemize}
    \item Công khai giá và các khoản phí trước khi khách hàng xác nhận mua.
    \item Không tự ý bổ sung sản phẩm vào đơn hàng.
    \item Không cấp quyền học chỉ dựa trên dữ liệu do Frontend gửi lên.
    \item Mọi giá trị thanh toán phải được Backend tính lại.
    \item Lưu lại snapshot thương mại tại thời điểm tạo đơn hàng.
    \item Không xử lý cùng một giao dịch thành công nhiều lần.
    \item Cung cấp cơ chế tra cứu trạng thái đơn hàng.
    \item Xử lý khiếu nại và hoàn tiền theo quy trình minh bạch.
\end{itemize}

% =========================================================
\section{Niêm yết giá}

\subsection{Thông tin giá}

Trang bán hàng phải hiển thị phù hợp:

\begin{itemize}
    \item Giá hiện tại.
    \item Giá gốc nếu có khuyến mại.
    \item Đơn vị tiền tệ.
    \item Hình thức và thời hạn gói.
    \item Thuế hoặc khoản phí nếu thuộc trường hợp áp dụng.
    \item Điều kiện gia hạn, nâng cấp hoặc hoàn tiền.
\end{itemize}

Giá hiển thị không mặc nhiên bao gồm phí do ngân hàng,
tổ chức phát hành hoặc cổng thanh toán thu trực tiếp,
trừ khi AILMS công bố khác.

\subsection{Thay đổi giá}

AILMS có thể điều chỉnh giá của sản phẩm chưa được mua.

Việc thay đổi giá không làm thay đổi đơn hàng đã thanh toán thành công,
trừ trường hợp:

\begin{itemize}
    \item Có thỏa thuận với khách hàng.
    \item Cần sửa lỗi giao dịch theo quy định pháp luật.
    \item Có hành vi gian lận hoặc can thiệp trái phép.
\end{itemize}

\subsection{Lỗi hiển thị giá}

Nếu hệ thống hiển thị mức giá sai rõ ràng do lỗi kỹ thuật,
AILMS phải thông báo cho khách hàng và có thể:

\begin{itemize}
    \item Cho phép mua theo giá đúng sau khi khách hàng xác nhận.
    \item Hủy đơn hàng chưa thanh toán.
    \item Hoàn lại toàn bộ số tiền nếu giao dịch đã phát sinh
    nhưng không thể thực hiện hợp lý.
\end{itemize}

AILMS không nên tự ý thu thêm tiền sau khi giao dịch đã hoàn tất
nếu chưa có sự đồng ý hợp lệ của khách hàng.

% =========================================================
\section{Khuyến mại}

Chương trình khuyến mại phải công bố:

\begin{itemize}
    \item Tên chương trình.
    \item Đối tượng áp dụng.
    \item Sản phẩm hoặc gói được áp dụng.
    \item Mức ưu đãi.
    \item Thời gian bắt đầu và kết thúc.
    \item Giới hạn lượt sử dụng.
    \item Điều kiện loại trừ nếu có.
\end{itemize}

Việc hiển thị giá gốc và giá khuyến mại phải trung thực,
không tạo cảm giác giảm giá giả tạo hoặc gây nhầm lẫn.

Khuyến mại có thể kết thúc khi:

\begin{itemize}
    \item Hết thời hạn.
    \item Đạt số lượng sử dụng tối đa.
    \item Hết ngân sách hoặc số lượng ưu đãi đã công bố.
    \item Chương trình bị tạm dừng theo căn cứ hợp lệ.
\end{itemize}

% =========================================================
\section{Mã giảm giá}

\subsection{Loại mã giảm giá}

AILMS có thể hỗ trợ:

\begin{itemize}
    \item Giảm theo phần trăm.
    \item Giảm số tiền cố định.
    \item Giảm theo khóa học hoặc gói cụ thể.
    \item Giảm cho danh mục sản phẩm.
    \item Giảm cho nhóm người dùng đủ điều kiện.
    \item Mã sử dụng toàn hệ thống.
\end{itemize}

\subsection{Điều kiện sử dụng}

Một Coupon có thể được giới hạn bởi:

\begin{itemize}
    \item Thời gian hiệu lực.
    \item Số lượt sử dụng toàn hệ thống.
    \item Số lượt sử dụng trên mỗi tài khoản.
    \item Giá trị đơn hàng tối thiểu.
    \item Giá trị giảm tối đa.
    \item Danh sách sản phẩm được áp dụng.
    \item Vai trò hoặc nhóm khách hàng.
    \item Trạng thái tài khoản.
\end{itemize}

\subsection{Quy tắc áp dụng}

Trừ khi chương trình quy định khác:

\begin{itemize}
    \item Mỗi đơn hàng chỉ được áp dụng một Coupon.
    \item Coupon không được quy đổi thành tiền mặt.
    \item Giá trị giảm không vượt quá phần giá trị được phép giảm.
    \item Coupon không được chuyển nhượng nếu được cấp riêng cho tài khoản.
    \item Coupon hết hạn không được khôi phục mặc nhiên.
\end{itemize}

Backend phải kiểm tra lại toàn bộ điều kiện Coupon tại thời điểm tạo đơn hàng
và trước khi xác nhận thanh toán nếu cần thiết.

\subsection{Giữ và hoàn trả lượt sử dụng}

Nếu hệ thống giữ lượt Coupon khi tạo đơn:

\begin{itemize}
    \item Lượt sử dụng được đặt ở trạng thái tạm giữ.
    \item Khi thanh toán thành công, lượt sử dụng được xác nhận.
    \item Khi đơn hết hạn hoặc bị hủy hợp lệ,
    lượt sử dụng có thể được giải phóng.
    \item Coupon đã hết hạn trong thời gian đơn chờ
    không nhất thiết được khôi phục sau khi đơn bị hủy.
\end{itemize}

% =========================================================
\section{Giỏ hàng}

\subsection{Thêm sản phẩm}

Người dùng có thể thêm sản phẩm vào giỏ nếu:

\begin{itemize}
    \item Sản phẩm còn được bán.
    \item Gói đào tạo đang hoạt động.
    \item Người dùng chưa sở hữu quyền lợi tương đương còn hiệu lực.
    \item Lớp học còn khả năng tiếp nhận hoặc hỗ trợ Waitlist.
    \item Tài khoản đáp ứng điều kiện đăng ký.
\end{itemize}

\subsection{Giới hạn giỏ hàng}

Hệ thống có thể từ chối:

\begin{itemize}
    \item Hai gói xung đột của cùng một khóa học.
    \item Sản phẩm người dùng đã sở hữu.
    \item Sản phẩm đã ngừng bán.
    \item Sản phẩm không còn đáp ứng điều kiện.
    \item Giỏ hàng vượt giới hạn số lượng cho phép.
\end{itemize}

\subsection{Giá trong giỏ hàng}

Giá hiển thị trong giỏ hàng chỉ mang tính tạm thời
cho đến khi đơn hàng được tạo thành công.

Thay đổi về giá, Coupon, sĩ số hoặc trạng thái sản phẩm
có thể được cập nhật khi người dùng chuyển tới bước thanh toán.

% =========================================================
\section{Tạo đơn hàng}

\subsection{Snapshot thương mại}

Khi tạo đơn hàng, hệ thống phải lưu snapshot tối thiểu:

\begin{itemize}
    \item Mã và tên sản phẩm.
    \item Mã khóa học và gói đào tạo.
    \item Hình thức đào tạo.
    \item Giá gốc.
    \item Giá bán tại thời điểm đặt hàng.
    \item Giá trị Coupon.
    \item Thuế và phí nếu áp dụng.
    \item Tổng số tiền phải thanh toán.
    \item Đơn vị tiền tệ.
\end{itemize}

Snapshot giúp bảo đảm thay đổi giá sau này
không làm sai lệch đơn hàng đã tạo.

\subsection{Tính tổng đơn hàng}

Tổng thanh toán có thể được tính theo nguyên tắc:

\[
\text{Tổng thanh toán}
=
\text{Tổng giá bán}
-
\text{Giảm giá hợp lệ}
+
\text{Thuế và phí áp dụng}
\]

Backend là nguồn quyết định cuối cùng đối với kết quả tính toán.

Frontend không được tự gửi một mức tổng tiền để Backend tin cậy
mà không kiểm tra lại.

\subsection{Trạng thái đơn hàng}

Đơn hàng có thể có các trạng thái:

\begin{itemize}
    \item Chờ thanh toán.
    \item Đang xử lý.
    \item Đã thanh toán.
    \item Thanh toán thất bại.
    \item Đã hủy.
    \item Đã hết hạn.
    \item Hoàn tiền một phần.
    \item Đã hoàn tiền.
    \item Đang tranh chấp.
\end{itemize}

Tên enum thực tế phải thống nhất với mã nguồn và cơ sở dữ liệu.

% =========================================================
\section{Thời hạn thanh toán}

Đơn hàng chờ thanh toán có thể có thời hạn, ví dụ 15 phút.

Thời hạn cụ thể phải được hiển thị cho khách hàng
hoặc quy định trong cấu hình giao dịch.

Khi hết thời hạn, hệ thống có thể:

\begin{itemize}
    \item Đánh dấu đơn hàng hết hạn.
    \item Hủy phiên thanh toán chưa hoàn tất.
    \item Giải phóng Coupon đang giữ.
    \item Giải phóng chỗ giữ tạm trong lớp.
    \item Yêu cầu người dùng tạo đơn hàng mới.
\end{itemize}

Việc đơn hàng AILMS hết hạn không bảo đảm rằng giao dịch
tại cổng thanh toán đã tự động bị hủy.

Nếu nhận kết quả thanh toán muộn, Backend phải đối chiếu
trước khi cấp quyền hoặc hoàn trả tiền.

% =========================================================
\section{Giữ chỗ trong lớp học}

Việc thêm gói lớp nhóm vào giỏ hàng không mặc nhiên giữ chỗ.

Nếu hệ thống hỗ trợ giữ chỗ khi tạo đơn:

\begin{itemize}
    \item Chỗ chỉ được giữ trong thời hạn đơn hàng.
    \item Chỗ được xác nhận sau khi thanh toán thành công.
    \item Chỗ được giải phóng khi đơn hết hạn hoặc bị hủy.
    \item Hệ thống phải kiểm soát đồng thời để tránh vượt sĩ số.
\end{itemize}

Nếu thanh toán thành công nhưng không thể xác nhận chỗ
do lỗi đồng thời hoặc sai lệch dữ liệu, AILMS phải:

\begin{itemize}
    \item Đề xuất lớp tương đương.
    \item Đưa học viên vào Waitlist sau khi có sự đồng ý.
    \item Hoàn tiền nếu không cung cấp được phương án phù hợp.
\end{itemize}

% =========================================================
\section{Nâng cấp và gia hạn gói}

\subsection{Nâng cấp}

Học viên có thể được nâng cấp từ gói hiện tại sang gói có quyền lợi cao hơn,
chẳng hạn:

\begin{itemize}
    \item Từ tự học sang lớp nhóm.
    \item Từ tự học sang một-một.
    \item Từ lớp nhóm sang gói có thêm buổi tư vấn.
\end{itemize}

Mức chênh lệch không nên chỉ được tính bằng phép trừ đơn giản
nếu quyền lợi cũ đã được sử dụng một phần.

Hệ thống có thể xem xét:

\begin{itemize}
    \item Giá đã thanh toán thực tế.
    \item Coupon đã sử dụng.
    \item Thời gian hoặc nội dung đã sử dụng.
    \item Quyền lợi đã phát sinh.
    \item Mức giá hiện tại của gói mới.
    \item Thuế và điều kiện chương trình.
\end{itemize}

Công thức minh họa:

\[
\text{Số tiền nâng cấp}
=
\max\left(
0,\,
\text{Giá gói mới}
-
\text{Giá trị được ghi nhận từ gói cũ}
\right)
\]

Giá trị được ghi nhận từ gói cũ phải được quy định rõ,
không mặc nhiên bằng toàn bộ số tiền từng thanh toán.

\subsection{Gia hạn}

Gia hạn tạo thêm hoặc kéo dài thời hạn quyền truy cập.

Trước khi xác nhận, hệ thống phải hiển thị:

\begin{itemize}
    \item Thời gian gia hạn.
    \item Ngày hết hạn mới.
    \item Giá và ưu đãi.
    \item Thời điểm bắt đầu kỳ gia hạn.
\end{itemize}

AILMS không tự động gia hạn và thu tiền định kỳ
nếu người dùng chưa đồng ý rõ ràng với cơ chế thanh toán lặp lại.

% =========================================================
\section{Phương thức thanh toán}

AILMS có thể hỗ trợ:

\begin{itemize}
    \item Cổng thanh toán điện tử.
    \item Ví điện tử.
    \item Thẻ thanh toán.
    \item Chuyển khoản.
    \item Phương thức khác được công bố.
\end{itemize}

Khả năng sử dụng phụ thuộc vào:

\begin{itemize}
    \item Quốc gia và đơn vị tiền tệ.
    \item Nhà cung cấp thanh toán.
    \item Trạng thái tích hợp.
    \item Điều kiện của tài khoản người mua.
\end{itemize}

Thông tin thẻ hoặc tài khoản thanh toán nhạy cảm
có thể được xử lý trực tiếp bởi nhà cung cấp thanh toán.

% =========================================================
\section{Quy trình thanh toán PayPal}

Một giao dịch PayPal có thể diễn ra theo trình tự:

\begin{enumerate}
    \item Hệ thống tạo Order AILMS.
    \item Backend tạo PayPal Order.
    \item PayPal trả về Approval URL.
    \item Người dùng được chuyển đến PayPal.
    \item Người dùng đăng nhập và chấp thuận thanh toán.
    \item PayPal chuyển người dùng về Return URL.
    \item Frontend gửi PayPal Order ID tới Backend.
    \item Backend gọi Capture trực tiếp tới PayPal.
    \item Backend kiểm tra trạng thái và số tiền.
    \item Đơn hàng được đánh dấu đã thanh toán.
    \item Hệ thống cấp quyền học tương ứng.
\end{enumerate}

Việc người dùng quay lại Return URL không chứng minh thanh toán thành công.

Backend chỉ cấp quyền sau khi nhận và xác minh kết quả hợp lệ
từ nhà cung cấp thanh toán.

% =========================================================
\section{Approval URL và tạo lại phiên thanh toán}

Approval URL có thể hết hạn hoặc không còn hợp lệ.

Hệ thống chỉ nên tạo lại Approval URL khi:

\begin{itemize}
    \item Đơn hàng vẫn còn hiệu lực.
    \item Đơn hàng chưa được thanh toán.
    \item Không có phiên thanh toán đang được xử lý thành công.
    \item Người yêu cầu là chủ đơn hàng hoặc có quyền phù hợp.
\end{itemize}

Khi tạo PayPal Order mới cho cùng một Order AILMS:

\begin{itemize}
    \item Phiên cũ phải được đánh dấu không còn được ưu tiên.
    \item Hệ thống phải lưu liên kết giữa các lần thử thanh toán.
    \item Chỉ một kết quả thành công được dùng để hoàn tất đơn hàng.
\end{itemize}

% =========================================================
\section{Xác minh kết quả thanh toán}

Trước khi đánh dấu đơn hàng đã thanh toán, Backend phải kiểm tra:

\begin{itemize}
    \item Mã giao dịch của nhà cung cấp.
    \item Mã Order AILMS liên quan.
    \item Trạng thái giao dịch.
    \item Số tiền và đơn vị tiền tệ.
    \item Tài khoản merchant nhận tiền.
    \item Giao dịch đã từng được xử lý hay chưa.
\end{itemize}

Không tin cậy trực tiếp:

\begin{itemize}
    \item Trạng thái do Frontend truyền lên.
    \item Tham số trên Return URL.
    \item Ảnh chụp giao dịch.
    \item Nội dung phản hồi do người dùng tự nhập.
\end{itemize}

% =========================================================
\section{Idempotency và chống xử lý trùng}

Mỗi yêu cầu thanh toán quan trọng phải có mã định danh duy nhất.

Hệ thống phải bảo đảm:

\begin{itemize}
    \item Một Payment ID chỉ được hoàn tất một lần.
    \item Một callback gửi lại không cấp quyền học lần thứ hai.
    \item Capture lặp lại không tạo nhiều Enrollment.
    \item Hoàn tiền lặp lại không vượt quá số tiền có thể hoàn.
    \item Mỗi mã giao dịch bên ngoài liên kết duy nhất với giao dịch nội bộ phù hợp.
\end{itemize}

Khi nhận yêu cầu trùng, hệ thống nên trả lại kết quả đã xử lý
hoặc trạng thái hiện tại thay vì thực thi lại toàn bộ nghiệp vụ.

% =========================================================
\section{Giao dịch Sandbox}

Môi trường Sandbox chỉ phục vụ:

\begin{itemize}
    \item Kiểm thử luồng thanh toán.
    \item Kiểm thử approval và capture.
    \item Kiểm thử giao dịch thất bại.
    \item Kiểm thử hoàn tiền.
    \item Kiểm thử idempotency và lỗi mạng.
\end{itemize}

Dữ liệu Sandbox:

\begin{itemize}
    \item Không có giá trị tiền thật.
    \item Không được dùng làm chứng từ tài chính.
    \item Không cấp quyền trên môi trường Production.
    \item Phải sử dụng credentials tách biệt với Production.
\end{itemize}

% =========================================================
\section{Cấp quyền sau thanh toán}

Sau khi giao dịch được xác nhận:

\begin{itemize}
    \item Đơn hàng được cập nhật trạng thái đã thanh toán.
    \item Payment được lưu trạng thái hoàn tất.
    \item Coupon được xác nhận đã sử dụng.
    \item Sản phẩm đã mua được xóa khỏi giỏ hàng.
    \item Enrollment hoặc quyền sở hữu được tạo.
    \item Thông báo được gửi cho học viên.
\end{itemize}

Theo hình thức đào tạo:

\begin{itemize}
    \item \textbf{Tự học:} Cấp quyền truy cập nội dung.
    \item \textbf{Lớp nhóm:} Thêm học viên vào lớp và thông báo cho các bên.
    \item \textbf{Một-một:} Tạo yêu cầu xếp lịch hoặc hồ sơ tư vấn.
\end{itemize}

Các bước cấp quyền phải có khả năng chạy lại an toàn
nếu một phần quy trình bị gián đoạn.

% =========================================================
\section{Hóa đơn, biên lai và xác nhận thanh toán}

\subsection{Xác nhận giao dịch}

Sau khi thanh toán thành công, hệ thống có thể cung cấp:

\begin{itemize}
    \item Mã đơn hàng.
    \item Mã giao dịch.
    \item Thời điểm thanh toán.
    \item Sản phẩm đã mua.
    \item Số tiền và đơn vị tiền tệ.
    \item Phương thức thanh toán.
    \item Trạng thái giao dịch.
\end{itemize}

\subsection{Biên lai điện tử}

Biên lai hoặc xác nhận dạng PDF do hệ thống tạo
chỉ phản ánh dữ liệu giao dịch trên AILMS.

Biên lai nội bộ không mặc nhiên là hóa đơn giá trị gia tăng
hoặc hóa đơn điện tử hợp pháp.

\subsection{Hóa đơn điện tử}

Hóa đơn điện tử chỉ được phát hành khi:

\begin{itemize}
    \item Đơn vị cung cấp dịch vụ có nghĩa vụ và đủ điều kiện phát hành.
    \item Thông tin người mua được cung cấp đúng hạn.
    \item Dữ liệu đáp ứng quy định về hóa đơn và thuế.
    \item Hệ thống hoặc nhà cung cấp hóa đơn điện tử đã được tích hợp hợp lệ.
\end{itemize}

Khách hàng có trách nhiệm cung cấp chính xác:

\begin{itemize}
    \item Tên cá nhân hoặc tổ chức.
    \item Mã số thuế nếu áp dụng.
    \item Địa chỉ và thông tin nhận hóa đơn.
    \item Thông tin khác theo yêu cầu hợp lệ.
\end{itemize}

% =========================================================
\section{Hủy đơn hàng}

\subsection{Đơn hàng chưa thanh toán}

Người dùng có thể hủy đơn chưa thanh toán nếu:

\begin{itemize}
    \item Đơn chưa được cổng thanh toán xác nhận.
    \item Không có giao dịch đang capture.
    \item Không có điều kiện đặc biệt ngăn việc hủy.
\end{itemize}

\subsection{Đơn hàng đã thanh toán}

Đơn đã thanh toán không được xử lý bằng thao tác hủy đơn thông thường.

Nếu khách hàng không muốn tiếp tục sử dụng,
phải thực hiện theo quy trình hoàn tiền hoặc hủy đăng ký.

% =========================================================
\section{Hoàn tiền}

\subsection{Điều kiện}

Yêu cầu hoàn tiền được xem xét dựa trên:

\begin{itemize}
    \item Thời gian kể từ ngày kích hoạt.
    \item Loại gói đào tạo.
    \item Tỷ lệ và quyền lợi đã sử dụng.
    \item Số buổi đã thực hiện.
    \item Điều kiện công bố tại thời điểm mua.
    \item Lý do và bằng chứng của yêu cầu.
\end{itemize}

\subsection{Hoàn toàn phần và một phần}

\begin{itemize}
    \item \textbf{Hoàn toàn phần:} Hoàn toàn bộ khoản đủ điều kiện.
    \item \textbf{Hoàn một phần:} Hoàn phần giá trị chưa sử dụng
    sau các khoản điều chỉnh hợp lệ.
\end{itemize}

Tổng số tiền hoàn không được vượt quá số tiền đã thu thành công
cho phần dịch vụ liên quan.

\subsection{Quy trình}

\begin{enumerate}
    \item Khách hàng gửi yêu cầu.
    \item Hệ thống kiểm tra đơn hàng và quyền sử dụng.
    \item Bộ phận có thẩm quyền phê duyệt hoặc từ chối.
    \item Backend gửi yêu cầu hoàn tới cổng thanh toán.
    \item Kết quả hoàn tiền được xác minh.
    \item Đơn hàng và Payment được cập nhật.
    \item Quyền học tương ứng được thu hồi hoặc điều chỉnh.
\end{enumerate}

Không đánh dấu đã hoàn tiền chỉ vì yêu cầu hoàn đã được gửi.
Phải phân biệt trạng thái đang xử lý và hoàn tất.

% =========================================================
\section{Chargeback và tranh chấp thanh toán}

\subsection{Tiếp nhận Chargeback}

Khi nhận thông báo Chargeback, AILMS có thể:

\begin{itemize}
    \item Đánh dấu giao dịch đang tranh chấp.
    \item Tạm hạn chế quyền lợi phát sinh từ giao dịch.
    \item Bảo toàn nhật ký và bằng chứng.
    \item Đối chiếu việc sử dụng khóa học.
    \item Phản hồi nhà cung cấp thanh toán trong thời hạn.
    \item Liên hệ khách hàng để xác minh.
\end{itemize}

\subsection{Không mặc định coi khách hàng gian lận}

Chargeback có thể phát sinh do:

\begin{itemize}
    \item Chủ tài khoản không nhận ra giao dịch.
    \item Tài khoản thanh toán bị sử dụng trái phép.
    \item Dịch vụ không được cung cấp đúng cam kết.
    \item Giao dịch bị trùng.
    \item Tranh chấp hợp pháp khác.
\end{itemize}

Vì vậy, AILMS không nên kết luận gian lận
chỉ dựa trên việc có yêu cầu Chargeback.

\subsection{Biện pháp xử lý}

Nếu Chargeback được xác định bất lợi cho giao dịch:

\begin{itemize}
    \item Quyền truy cập mua từ giao dịch có thể bị thu hồi.
    \item Enrollment liên quan có thể chuyển trạng thái đã hủy.
    \item Khoản Revenue Share chưa thanh toán có thể được điều chỉnh.
    \item Tài khoản có thể bị hạn chế nếu có bằng chứng lạm dụng.
\end{itemize}

Chứng chỉ không nên tự động bị thu hồi chỉ vì Chargeback.

Việc thu hồi chứng chỉ chỉ nên thực hiện khi:

\begin{itemize}
    \item Quyền học bị xác định là không hợp lệ.
    \item Có gian lận học tập hoặc thanh toán đã được xác minh.
    \item Chứng chỉ được cấp do lỗi hoặc dữ liệu không hợp lệ.
    \item Có căn cứ khác theo chính sách đào tạo.
\end{itemize}

% =========================================================
\section{Phát hiện giao dịch nghi ngờ}

Dấu hiệu rủi ro có thể gồm:

\begin{itemize}
    \item Nhiều lần thanh toán thất bại trong thời gian ngắn.
    \item Một phương thức thanh toán được dùng cho nhiều tài khoản bất thường.
    \item Giá trị giao dịch hoặc vị trí truy cập bất thường.
    \item Nhiều yêu cầu hoàn tiền hoặc Chargeback.
    \item Can thiệp vào số tiền hoặc callback.
    \item Dùng tài khoản chiếm đoạt để mua khóa học.
\end{itemize}

Khi phát hiện rủi ro, hệ thống có thể:

\begin{itemize}
    \item Yêu cầu xác minh bổ sung.
    \item Tạm giữ việc cấp quyền.
    \item Tạm khóa chức năng thanh toán.
    \item Chuyển giao dịch sang kiểm tra thủ công.
    \item Hủy giao dịch chưa hoàn tất nếu có căn cứ.
\end{itemize}

Biện pháp phải tương xứng với mức độ rủi ro
và cho phép người dùng hợp pháp giải trình.

% =========================================================
\section{Chia sẻ doanh thu}

\subsection{Điều kiện áp dụng}

Revenue Share chỉ áp dụng khi có:

\begin{itemize}
    \item Hợp đồng hoặc thỏa thuận hợp lệ.
    \item Khóa học thuộc phạm vi chia sẻ.
    \item Tỷ lệ hoặc công thức được xác định.
    \item Kỳ đối soát và điều kiện thanh toán rõ ràng.
\end{itemize}

\subsection{Cơ sở tính}

Doanh thu chia sẻ có thể được tính trên giá trị thực thu:

\[
\text{Doanh thu chia sẻ}
=
\text{Giá trị thực thu đủ điều kiện}
\times
\text{Tỷ lệ chia sẻ}
\]

Giá trị thực thu đủ điều kiện có thể được xác định từ:

\[
\begin{aligned}
\text{Giá trị thực thu đủ điều kiện}
=&\ \text{Số tiền thu thành công}\\
&-\text{Khoản hoàn tiền}\\
&-\text{Chargeback}\\
&-\text{Thuế, phí và khoản loại trừ theo hợp đồng}
\end{aligned}
\]

Khoản được loại trừ phải được ghi rõ trong hợp đồng,
không được tự ý phát sinh sau kỳ bán hàng.

\subsection{Phân bổ đơn hàng nhiều sản phẩm}

Nếu đơn hàng có nhiều sản phẩm và một Coupon dùng cho toàn đơn,
khoản giảm có thể được phân bổ theo:

\begin{itemize}
    \item Tỷ trọng giá bán của từng sản phẩm.
    \item Quy tắc ưu tiên của chương trình.
    \item Phương pháp khác được công bố trong thỏa thuận.
\end{itemize}

Phương pháp phải nhất quán và có khả năng kiểm tra lại.

% =========================================================
\section{Kỳ đối soát}

Kỳ đối soát có thể được thực hiện hàng tháng.

Báo cáo đối soát nên bao gồm:

\begin{itemize}
    \item Danh sách đơn hàng thành công.
    \item Khóa học và gói liên quan.
    \item Giá bán và giảm giá.
    \item Số tiền thực thu.
    \item Hoàn tiền và Chargeback.
    \item Tỷ lệ chia sẻ.
    \item Số tiền dự kiến thanh toán.
    \item Khoản điều chỉnh từ kỳ trước.
\end{itemize}

Giảng viên hoặc đối tác phải có thời hạn hợp lý để:

\begin{itemize}
    \item Xem báo cáo.
    \item Gửi yêu cầu đối soát.
    \item Cung cấp bằng chứng.
    \item Nhận kết quả điều chỉnh.
\end{itemize}

% =========================================================
\section{Thanh toán hoa hồng}

Khoản chia sẻ doanh thu có thể:

\begin{itemize}
    \item Được chuyển cùng kỳ lương.
    \item Được thanh toán riêng theo hợp đồng.
    \item Được chuyển sang kỳ sau nếu chưa đạt ngưỡng tối thiểu.
\end{itemize}

Việc gộp với kỳ lương chỉ phù hợp khi giảng viên
đồng thời là nhân sự có quan hệ lao động và hợp đồng có quy định.

Đối với cộng tác viên hoặc đối tác,
cần sử dụng cơ chế thanh toán phù hợp với bản chất hợp đồng.

Các nghĩa vụ thuế, chứng từ và khấu trừ
được thực hiện theo pháp luật và thỏa thuận áp dụng.

% =========================================================
\section{Bảo mật dữ liệu thanh toán}

AILMS phải:

\begin{itemize}
    \item Không ghi secret key vào mã nguồn.
    \item Không gửi secret key tới Frontend.
    \item Tách credentials Sandbox và Production.
    \item Xác minh webhook hoặc callback.
    \item Giới hạn quyền truy cập dữ liệu tài chính.
    \item Che bớt thông tin nhạy cảm khi hiển thị.
    \item Ghi Audit Log cho các thao tác tài chính quan trọng.
\end{itemize}

Không ghi vào log:

\begin{itemize}
    \item Mật khẩu.
    \item Secret key.
    \item Access Token đầy đủ.
    \item Thông tin thẻ đầy đủ.
    \item Dữ liệu xác thực không cần thiết.
\end{itemize}

% =========================================================
\section{Audit Log tài chính}

Audit Log nên ghi nhận:

\begin{itemize}
    \item Tạo và hủy đơn hàng.
    \item Áp dụng hoặc giải phóng Coupon.
    \item Tạo phiên thanh toán.
    \item Capture thành công hoặc thất bại.
    \item Phê duyệt và thực hiện hoàn tiền.
    \item Điều chỉnh trạng thái thanh toán.
    \item Đối soát và phê duyệt Revenue Share.
    \item Xử lý Chargeback.
\end{itemize}

Nhật ký phải thể hiện:

\begin{itemize}
    \item Người hoặc hệ thống thực hiện.
    \item Thời điểm.
    \item Đối tượng giao dịch.
    \item Giá trị trước và sau.
    \item Mã tham chiếu.
    \item Lý do điều chỉnh nếu có.
\end{itemize}

% =========================================================
\section{Khiếu nại giao dịch}

Khách hàng có thể khiếu nại về:

\begin{itemize}
    \item Thanh toán bị trừ tiền nhưng chưa được cấp quyền.
    \item Thanh toán trùng.
    \item Sai giá hoặc sai Coupon.
    \item Không nhận được sản phẩm.
    \item Hoàn tiền chậm hoặc không chính xác.
    \item Hóa đơn hoặc biên lai có sai sót.
\end{itemize}

Yêu cầu nên kèm theo:

\begin{itemize}
    \item Mã đơn hàng.
    \item Mã giao dịch.
    \item Thời điểm thanh toán.
    \item Số tiền.
    \item Mô tả sự việc.
    \item Bằng chứng phù hợp.
\end{itemize}

AILMS hướng tới xác nhận tiếp nhận trong vòng 48 giờ làm việc.
Thời gian giải quyết cuối cùng phụ thuộc vào cổng thanh toán,
ngân hàng và mức độ phức tạp.

% =========================================================
\section{Sửa đổi chính sách}

Chính sách có thể được cập nhật khi:

\begin{itemize}
    \item Thay đổi phương thức thanh toán.
    \item Thay đổi mô hình giá hoặc khuyến mại.
    \item Bổ sung gia hạn hoặc thanh toán định kỳ.
    \item Thay đổi quy trình hóa đơn.
    \item Có yêu cầu pháp luật mới.
\end{itemize}

Thay đổi quan trọng ảnh hưởng đến giao dịch tương lai
phải được công bố trước khi khách hàng xác nhận mua.

Quyền lợi của đơn hàng đã hoàn thành không nên bị thay đổi hồi tố,
trừ khi có căn cứ pháp lý hoặc thỏa thuận hợp lệ.

% =========================================================
\section{Thông tin liên hệ}

Khách hàng và đối tác có thể liên hệ:

\begin{itemize}
    \item \textbf{Đơn vị phụ trách:}
    Phòng Tài chính và Kinh doanh AILMS.

    \item \textbf{Email thanh toán:}
    \href{mailto:billing@ailms.example}{billing@ailms.example}.

    \item \textbf{Email đối soát:}
    \href{mailto:reconciliation@ailms.example}
    {reconciliation@ailms.example}.

    \item \textbf{Kênh hỗ trợ:}
    Chức năng gửi yêu cầu trên hệ thống.
\end{itemize}

\textcolor{ailmsred}{
\textit{Lưu ý: Thay địa chỉ thư điện tử mẫu, tên pháp nhân,
thông tin thuế và thời hạn nghiệp vụ bằng dữ liệu chính thức
trước khi công bố.}
}

% =========================================================
\section{Hiệu lực thi hành}

Chính sách có hiệu lực kể từ ngày được phê duyệt và công bố.

Khách hàng, giảng viên, đối tác và bộ phận liên quan
có trách nhiệm tuân thủ chính sách trong phạm vi áp dụng.

\vfill

\begin{center}
    \textbf{ĐẠI DIỆN PHÒNG TÀI CHÍNH VÀ KINH DOANH AILMS}\\[1.5cm]
    \textit{Ký và ghi rõ họ tên}
\end{center}

\end{document}